%! Author = sriadityavedantam
%! Date = 3/25/26

\section{Homework 0}

\begin{problem}{
    \label{ex:0.1}
    The transitive property for the rationals holds: if $(a,b)\sim(\alpha,\beta)$ and $(\alpha,\beta)\sim(c,d)$, then $(a,b)\sim(c,d)$.
}
    Assume
    \[
        (a,b)\sim(\alpha,\beta)
        \qquad\text{and}\qquad
        (\alpha,\beta)\sim(c,d).
    \]
    By definition,
    \[
        a\beta=\alpha b
        \qquad\text{and}\qquad
        \alpha d=\beta c.
    \]
    Multiply the two equations:
    \[
        (a\beta)(\alpha d)=(\alpha b)(\beta c).
    \]
    Cancel the common factors:
    \[
        (a\cancel{\beta})(\cancel{\alpha}d)=(\cancel{\alpha}b)(\cancel{\beta}c).
    \]
    Hence
    \[
        ad=bc.
    \]
    Therefore
    \[
        (a,b)\sim(c,d).
    \]
\end{problem}

\begin{problem}[Lemma 1.6 --- Transitivity Proof]{
    The transitive property for the rationals holds if $(a,b)\sim(\alpha,\beta)$ and $(\alpha,\beta)\sim(c,d)$, then $(a,b)\sim(c,d)$.
}
    Assume
    \[
        a\beta=\alpha b
        \qquad\text{and}\qquad
        \alpha d=\beta c.
    \]
    Then
    \begin{align*}
        (a\beta)(\alpha d)
&=(\alpha b)(\beta c)\\
(a\cancel{\beta})(\cancel{\alpha}d)
&=(\cancel{\alpha}b)(\cancel{\beta}c)\\
ad&=bc.
    \end{align*}
    Hence $(a,b)\sim(c,d)$.
\end{problem}

\begin{problem}[Proposition 1.8 --- Addition is Well-Defined]{
    Given $(a,b)\sim(c,d)$ and $(\alpha,\beta)\sim(\gamma,\delta)$, show
    \[
        (a\beta+\alpha b,\,b\beta)\sim(c\delta+\gamma d,\,d\delta).
    \]
    This is how addition is defined:
    \[
        (a,b)+(\alpha,\beta)\coloneqq (a\beta+\alpha b,\,b\beta).
    \]
}
    From the relations,
    \[
        ad=bc
        \qquad\text{and}\qquad
        \alpha\delta=\beta\gamma.
    \]
    We compute
    \begin{align*}
        (a\beta+\alpha b)(d\delta)
&=a\beta d\delta+\alpha bd\delta\\
&=(ad)(\beta\delta)+(\alpha\delta)(bd).
    \end{align*}
    Using $ad=bc$ and $\alpha\delta=\beta\gamma$, this becomes
    \begin{align*}
        (ad)(\beta\delta)+(\alpha\delta)(bd)
    &=(bc)(\beta\delta)+(\beta\gamma)(bd)\\
        &=b\beta(c\delta+\gamma d).
    \end{align*}
    Therefore
    \[
        (a\beta+\alpha b)(d\delta)=(c\delta+\gamma d)(b\beta),
    \]
    so
    \[
        (a\beta+\alpha b,\,b\beta)\sim(c\delta+\gamma d,\,d\delta).
    \]
    Hence addition is well-defined.
\end{problem}

\begin{problem}{
    Why must we use $\z\times(\z\setminus\{0\})$?
    Show that the relation
    \[
        (a,b)\sim(c,d)\iff ad=bc
    \]
    is not an equivalence relation on $\z\times\z$.
}
    Suppose this relation is defined on $\z\times\z$.
    Take
    \[
        (a,b)=(1,0),\qquad (c,d)=(1,0),\qquad (e,f)=(1,1).
    \]
    Then
    \[
        (1,0)\sim(1,0)
    \]
    because
    \[
        1\cdot 0=0\cdot 1.
    \]
    Also
    \[
        (1,0)\sim(1,1)
    \]
    because
    \[
        1\cdot 1=0\cdot 1
    \]
    is false, so this choice does not work.
    Instead take
    \[
        (a,b)=(1,0),\qquad (c,d)=(0,0),\qquad (e,f)=(0,1).
    \]
    Then
    \[
        (1,0)\sim(0,0)
    \]
    because
    \[
        1\cdot 0=0\cdot 0.
    \]
    Also
    \[
        (0,0)\sim(0,1)
    \]
    because
    \[
        0\cdot 1=0\cdot 0.
    \]
    But
    \[
        (1,0)\not\sim(0,1)
    \]
    because
    \[
        1\cdot 1\neq 0\cdot 0.
    \]
    So transitivity fails.
    Hence this is not an equivalence relation on $\z\times\z$.
    This is why the second component must be nonzero.
\end{problem}